\section{Related work}
\label{sec:related}

\paragraph{Optimal market making.}
The stochastic-control literature following \citet{ho1981optimal} and
\citet{avellaneda2008high} models the arrival of fills to a quote at distance
$\delta$ from the mid-price as a Poisson process with intensity
$\lambda(\delta)=A e^{-k\delta}$. \citet{gueant2013dealing} solve the resulting
inventory problem in closed form; \citet{cartea2015algorithmic} give a
textbook treatment. In this framework the fill rule \emph{is} the intensity
function, and the queue in front of a quote does not appear at all.

\paragraph{Queues and fill probability.}
Queue-reactive and queueing models of the order book
\citep{cont2010stochastic,huang2015simulating} describe the dynamics of the
volume at each level; \citet{moallemi2016model} value a position in the queue.
\citet{lo2002econometric} model limit-order execution times with survival
analysis, which our outcome table is built to support. Queue imbalance at the
best prices predicts the next mid-price move \citep{gould2016queue}.
\todo{Add Zhong et al.\ (fill probability, queue position as a feature) once the
exact reference is confirmed.}

\paragraph{Backtesting practice.}
Retail and professional backtesting platforms fill limit orders on candle
touch and market orders at the next bar's open, with a user-chosen constant
slippage \todo{cite platform documentation: TradingView, backtrader,
vectorbt}. Event-driven simulators typically wait for traded volume at the
order's price to exceed the volume displayed ahead of it. Neither rule has, to
our knowledge, been validated against the outcomes of real orders.

\paragraph{Order-level data.}
Exchange order-level data has so far been available only under licence and
without participant identities. The \emph{Open Book} dataset
\citep{albers2026openbook} is, to our knowledge, the first public order-level
record of a major venue that follows every order of every participant.
